\begin{table}[htbp]
\centering
\small
\caption{Scenario-family sensitivity of the M4 top-decile set at the 10\% budget.
Each family varies one group of assumptions while holding the rest at the baseline
configuration. The table reports the minimum and maximum Jaccard overlap with the
baseline top decile across scenarios in the family.}
\label{tab:sensitivity}
\begin{threeparttable}
\begin{tabular}{@{}lcrr@{}}
\toprule
Family & Scenarios & Minimum $J$ & Maximum $J$ \\
\midrule
Aggregation (hybrid vs fractional) & 1 & 0.881 & 0.881 \\
Park threshold (500 m, 1,000 m, caps) & 4 & 0.521 & 0.682 \\
Weights (including equal weight) & 9 & 0.868 & 0.954 \\
QA exclusions (snap, coverage) & 3 & 0.969 & 0.996 \\
Missing park distance (median, cap) & 2 & 0.906 & 0.906 \\
Joint factorial (construct, cap, normalisation, aggregation) & 16 & 0.583 & 0.728 \\
Missing support & 1 & 1.000 & 1.000 \\
Normalisation (min--max, rank) & 3 & 0.734 & 0.814 \\
Population definition (count vs share) & 2 & 0.347 & 0.977 \\
Route policy (baseline, conservative, positive tags) & 5 & 0.420 & 1.000 \\
Thermal support (0.50, 0.25, strict) & 2 & 1.000 & 1.000 \\
Land cover (tree probability) & 1 & 0.782 & 0.782 \\
\bottomrule
\end{tabular}
\begin{tablenotes}[flushleft]\footnotesize
\item Jaccard values are computed against the baseline M4 top decile with the scenario's eligible set (eligibility-changing comparison). The joint factorial does not include route-policy variants, so the table does not support an all-family robustness claim. Two further families (entrance/barrier and thermal sampling) could not be evaluated because the required evidence was unavailable, and they are not reported as estimates.
\end{tablenotes}
\end{threeparttable}
\end{table}
